% CV - Darvin Cotrina Cervera
% Formato basado en las guías de la Office of Career Services (OCS), Harvard University:
% secciones Education -> Experience -> Skills, orden cronológico inverso,
% bullets con verbo de acción + resultado medible, sin pronombres, una página.
\documentclass[10pt,letterpaper]{article}

\usepackage[margin=0.55in]{geometry}
\usepackage[T1]{fontenc}
\usepackage{mathptmx}        % Times New Roman (fuente recomendada por Harvard OCS)
\usepackage{enumitem}
\usepackage{titlesec}
\usepackage[hidelinks]{hyperref}

\pagestyle{empty}
\setlength{\parindent}{0pt}

% Estilo de sección: NOMBRE EN VERSALITAS con línea horizontal debajo (estilo OCS)
\titleformat{\section}{\large\bfseries\scshape}{}{0pt}{}[\vspace{-4pt}\rule{\textwidth}{0.6pt}]
\titlespacing*{\section}{0pt}{6pt}{3pt}

% Cargo + fechas en la misma línea
\newcommand{\role}[2]{\textbf{#1}\hfill\textit{#2}\\[-2pt]}
% Empresa + ubicación
\newcommand{\employer}[2]{\textbf{#1}\hfill #2\\[-1pt]}

\begin{document}

% ---------- Encabezado ----------
\begin{center}
  {\LARGE\bfseries Darvin Cotrina Cervera}\\[2pt]
  {\small Gerente de Transformación Digital}\\[3pt]
  {\small Lima, Perú \textbar\ +51 948 269 619 \textbar\ \href{mailto:cc.darvin@gmail.com}{cc.darvin@gmail.com} \textbar\ \href{https://www.linkedin.com/in/darvin-cotrina}{linkedin.com/in/darvin-cotrina} \textbar\ \href{https://www.darvincotrina.com}{darvincotrina.com}}
\end{center}
\vspace{-8pt}

% ---------- Educación ----------
\section{Educación}
\textbf{MIT Professional Education} \hfill 2024--2025\\
Certificado Profesional en Transformación Digital\\[1pt]
\textbf{MITx on edX} \hfill 2020--2021\\
MicroMasters\textregistered\ Program in Statistics and Data Science\\[1pt]
\textbf{Universidad de Lima} \hfill 2004--2009\\
Ingeniería de Sistemas

% ---------- Experiencia ----------
\section{Experiencia}
\employer{TDP CORP S.A.}{Lima, Perú}
\role{Gerente de Transformación Digital}{oct 2025 -- Presente}
\begin{itemize}[leftmargin=1.4em, nosep]
  \item Dirigí la construcción de una plataforma digital propia con arquitectura AI-first y 7 aplicaciones integradas, desarrollada con agentes de programación (Claude, Codex, Kimi)
  \item Puse en producción 3 agentes de IA para reconciliación bancaria y facturación de compras y ventas, alcanzando 73\% de autonomía medible
  \item Lidero el despliegue de coworkers (agentes personales de IA) en las áreas de la organización para análisis de datos y asistencia operativa
  \item Diseñé el Plan Estratégico de Transformación Digital 2027 con horizonte 2030, aprobado por la Alta Dirección
\end{itemize}
\vspace{2pt}
\role{Jefe de Desarrollo e Innovación}{ago 2014 -- oct 2025}
\begin{itemize}[leftmargin=1.4em, nosep]
  \item Creé el área de Transformación Digital, formalizada y validada por la Alta Dirección como eje estratégico
  \item Lideré la migración de los sistemas de gestión a SAP Business One y la integración de ERP, CRM Zoho e inteligencia artificial
  \item Introduje la facturación electrónica, reduciendo más del 60\% los tiempos operativos
  \item Automaticé procesos con RPA, incluyendo la carga y generación de facturas electrónicas hacia SAP
  \item Dirigí el e-commerce corporativo, superando los 2,500 contactos generados desde la tienda en línea
\end{itemize}
\vspace{2pt}
\role{Analista de Sistemas}{mar 2012 -- sep 2014}
\begin{itemize}[leftmargin=1.4em, nosep]
  \item Coordiné la migración corporativa a G Suite y administré los sistemas de gestión internos
\end{itemize}
\vspace{2pt}
\role{Analista Programador}{ago 2011 -- mar 2012}
\begin{itemize}[leftmargin=1.4em, nosep]
  \item Desarrollé un Punto de Venta y un Sistema de Gestión Interno con generación de reportes
\end{itemize}

\vspace{3pt}
\employer{PRODUCTOS AVON S.A.}{Lima, Perú}
\role{Programador de Software}{sep 2008 -- abr 2009}
\begin{itemize}[leftmargin=1.4em, nosep]
  \item Diseñé un sistema de cruce de líneas telefónicas que redujo el proceso de 2 días a menos de 10 minutos
\end{itemize}

\vspace{3pt}
\employer{UNIVERSIDAD DE LIMA}{Lima, Perú}
\role{Soporte Técnico}{ene 2008 -- sep 2008}
\begin{itemize}[leftmargin=1.4em, nosep]
  \item Desarrollé utilitarios en C\#, VB y Java para los laboratorios académicos
\end{itemize}

% ---------- Certificaciones ----------
\section{Certificaciones}
\textbf{deeplearning.ai} --- Especialización en Inteligencia Artificial (2023); Deep Learning Specialization (2019--2020); Machine Learning Specialization\\[1pt]
\textbf{MIT Professional Education} --- Cultural Awareness: Relaciones Interculturales en Negocios Globales (2024)

% ---------- Habilidades ----------
\section{Habilidades}
\textbf{Estrategia y liderazgo:} Transformación digital, gestión del cambio, gobernanza de IA, PMO digital, metodologías ágiles (Scrum, DevOps), marcos COBIT, ITIL e ISO 27001\\[1pt]
\textbf{Tecnologías:} Agentes de IA, IA generativa multi-modelo, RPA, ERP (SAP Business One), CRM (Zoho), e-commerce, WMS, cloud computing, analítica de datos, Python, SQL, Git

\end{document}
